%% discussion_ko.tex
\chapter{논의}
\label{chap:discussion_ko}

\section{하이브리드 검색이 BM25와 Dense를 능가하는 이유}
\label{sec:disc_hybrid_ko}

모든 실험에 걸쳐 가장 일관된 발견은 Hybrid RRF가 두 단일 모달 검색 방법 모두를 능가한다는 것이다. Hybrid는 Recall@10$={0.745}$, MRR$={0.530}$, nDCG@10$={0.566}$을 달성하여 BM25 대비 각각 +4.5\%, +15.8\%, +11.3\%, Dense 대비 +0.6\%, +5.6\%, +3.4\% 향상하였다. 이 결과는 BM25와 밀집 검색이 경쟁적이 아닌 상보적 방법이라는 융합의 이론적 동기와 일치한다 \citep{cormack2009rrf}.

상보성은 난이도 요인 수준에서 가장 잘 드러난다. BM25는 정확한 식별자를 포함하는 질의에서 우수한 성능을 보인다: 문서 버전 충돌 질의에서 Recall@1$={0.583}$을 기록하며, 버전 번호와 정책 코드에 대한 키워드 정밀 매칭에서 이점이 있다. Dense는 개념적 관계의 의미적 이해가 필요한 질의에서 우수하다: 비교 질의에서 MRR$={0.730}$(BM25의 $0.618$ 대비). Hybrid는 RRF 융합이 두 시스템의 순위 정보를 보존하여 어느 신호에도 전적으로 의존하지 않기 때문에 모든 난이도 요인에서 최선 또는 근접 최선의 성능을 일관되게 달성한다.

RRF 융합 이점은 BM25와 Dense 모두 다른 이유로 정보를 제공하는 질의에서 가장 크게 나타난다. SEKD 질의 집합의 54%를 구성하는 다중 홉 및 표 종속성 시나리오가 이에 해당한다. 이러한 질의에서 BM25는 정확한 식별자(프로젝트 코드, 정책 섹션 참조)를 포함하는 문서를 검색하는 반면, Dense는 내용 질문과 의미적으로 일치하는 문서를 검색하며, RRF는 두 신호를 단일 순위 리스트로 병합한다.

하이브리드가 해결하지 못하는 유일한 실패 유형은 시간적 추론이다. Dense는 시간적 질의에서 Recall@10$={0.000}$을 달성하고, Hybrid의 $0.100$은 융합을 통한 미미한 BM25 기여를 반영한다. 시간적 질의는 384차원 밀집 임베딩에서 거의 보이지 않는 특정 날짜 토큰 매칭이 필요하다. 이는 시간적 내용(회의록, 버전 이력, 정책 개정 로그)이 상당히 많은 기업 코퍼스에서 BM25 구성 요소가 선택 사항이 아닌 필수 요소임을 시사한다.

\section{규칙 기반 계획기가 좋은 성능을 보인 이유}
\label{sec:disc_rule_ko}

7개 규칙 계획기는 SSA$={78.4\%}$와 Recall@10$={0.755}$를 달성하여 해당 지표에서 최고 성능을 기록하였다. 이러한 성공은 세 가지 상호 강화 요인에 기인한다.

\textbf{요인 1: Oracle 레이블이 규칙 논리를 반영한다.} Oracle 전략 레이블은 규칙 계획기의 설계를 그대로 반영하는 휴리스틱 과정을 통해 할당되었다. 구체적으로, oracle 할당에 사용된 추론 유형 매핑(시간적→BM25, 예외→Dense, 다중 홉→Hybrid)이 규칙 계획기의 규칙 R01--R04에 직접 인코딩되었다. 이 구조적 정렬은 규칙 계획기가 자체 설계 가정을 반영하는 정답에 대해 평가됨을 의미한다. 이는 oracle 구성 과정의 한계이지만(타당성 위협 참조), 추론 유형 분류체계가 잘 구조화된 기업 코퍼스에서 전략 선택에 충분한 신호를 제공한다는 진정한 통찰을 반영하기도 한다.

\textbf{요인 2: 높은 $n$을 가진 그룹의 완벽한 커버리지.} 규칙 계획기는 집계($n{=}35$), 예외($n{=}34$), 다중 홉($n{=}53$) 질의에서 100\% SSA를 달성한다. 이 세 그룹은 질의 집합의 32.1\%를 차지한다. 이러한 그룹에서 검색 전략은 명확하다: 집계와 다중 홉은 문서 전반의 커버리지가 필요하고(hybrid), 예외 질의는 의미적 조항 해석이 필요하다(dense). 결정론적 매핑은 설계상 이러한 경우에 정확하다.

\textbf{요인 3: 안전한 기본값으로서의 Hybrid.} 규칙 계획기의 기본 전략(R07)은 최선의 고정 기준선인 Hybrid이다. 규칙이 실행되지 않는 질의에 대해 Hybrid로 기본 설정함으로써, 규칙 계획기는 해당 질의에서 Hybrid보다 나쁜 성능을 보이지 않는다. Oracle 레이블의 65.5\%가 Hybrid이므로, 기본값이 다수 사례를 포착하고 전문 규칙이 고가치 예외를 처리한다.

계획기의 주요 약점은 BM25 과소 예측이다: oracle-BM25 질의 62개 중 5개만 BM25로 라우팅한다(8.1\% BM25 리콜). 5개의 시간적 질의는 올바르게 라우팅되지만, 강한 키워드 신호를 가진 범주(특정 API 문서, 정확한 숫자 조회)의 단일 홉 질의는 BM25가 경험적으로 이점이 있음에도 불구하고 Hybrid로 잘못 기본 설정된다. 키워드 밀집 질의 패턴을 대상으로 하는 추가 규칙을 통해 SSA를 80\% 이상으로 높일 수 있다.

\section{GPT-5가 규칙 계획기를 능가하지 못한 이유}
\label{sec:disc_gpt5_ko}

이 연구의 이론적으로 가장 중요한 발견은 GPT-5가 대규모 제로샷 추론 능력과 \$6.18의 추론 비용에도 불구하고 전체 검색 성능에서 7개 규칙 결정론적 시스템을 능가하지 못한다는 것이다. 이를 이해하기 위해서는 여러 혼재 요인을 분리해야 한다.

\textbf{과제 복잡성 불일치.} 추론 유형, 난이도 요인, 범주 등 구조화된 메타데이터를 기반으로 3개 범주에 대한 검색 전략 분류는 저복잡성 의사결정 과제이다. 구조화 프롬프트가 올바른 결정 논리를 명시적으로 인코딩하므로, LLM의 한계적 기여(자연어 이해, 맥락적 추론, 교차 도메인 추론)는 대부분의 질의에서 거의 무관하다. 50개의 레이블 예시로 훈련된 간단한 분류기가 무시할 수 있는 비용으로 비슷한 정확도를 달성할 가능성이 높다 \citep{shi2023large}.

\textbf{SSA가 휴리스틱 oracle 대비 정확성을 측정한다.} Oracle 레이블은 규칙 계획기와 동일한 논리를 반영한다. API 문서 질의에서 dense보다 bm25를 선호하는 GPT-5의 이탈은 계획 오류보다 휴리스틱에 대한 진정한 불동의를 나타낼 수 있다. GPT-5 오류의 44.7\%가 \texttt{oracle\_disagreement}로 분류된 것은 이 해석을 지지한다.

\textbf{GPT-5의 순위 품질 이점이 Recall@10 상한 효과에 의해 가려진다.} GPT-5는 MRR(+2.0\%)과 nDCG@10(+1.2\%), Recall@1(+2.9\%)과 Recall@5(+3.1\%)에서 규칙 계획기를 능가한다. 이러한 순위 품질 향상은 이미 두 시스템에서 높은 Recall@10(0.750과 0.755, 무시할 수 있는 0.7 pp 차이)보다 사용자 대면 응용에서 더 중요하다.

\textbf{시간적 질의 실패는 체계적인 보정 오류이다.} GPT-5는 시간적 질의에서 SSA 0\%를 달성하며, 일관되게 bm25 대신 hybrid를 예측한다. 구조화 프롬프트에는 명시적인 시간적 지침(``시간적→bm25'')이 포함되어 있지만, GPT-5의 추론 과정은 가장 일반적으로 옳은 기본값(hybrid, oracle의 65.5\%)을 선호하는 것으로 보인다. 이는 추론 모델에서 중요한 라우팅 규칙이 표준 프롬프트 지침보다 더 강하게 강화될 필요가 있음을 시사한다.

\textbf{API 문서 혼동이 지배적 오류 원인이다.} GPT-5 오류 85개 중 32개(37.6\%)가 API 문서 질의에서 발생한다. Oracle은 특정 질의 특성에 따라 세 가지 전략을 모두 할당하지만(정확한 엔드포인트 경로→bm25, 의미적 인증 질의→dense, 다중 홉 API 종속성→hybrid), GPT-5는 일관성 없이 전략을 선택한다.

\section{기업 RAG 시스템에 대한 시사점}
\label{sec:disc_implications_ko}

\textbf{시사점 1: Hybrid RRF가 올바른 기본 검색 아키텍처이다.} 이 연구에서 어떤 시스템도 BM25나 Dense에 단독으로 전념하는 것에서 이점을 얻지 못한다. 질의 분포가 알려지지 않은 기업 RAG 배포에서 Hybrid RRF는 두 검색 색인 유지의 설정 오버헤드 외에 추가 비용 없이 최선의 커버리지와 순위 품질 조합을 제공한다 \citep{luan2021sparse_dense}.

\textbf{시사점 2: 경량 규칙 계획기가 운영 비용 없이 가치를 추가한다.} 7개 규칙 계획기는 적당한 SSA(78.4\%)가 최선의 고정 기준선을 능가하기에 충분함을 보여준다(Recall@10: +1.4\%, Hit@10: +1.3\%). 규칙은 해석 가능하고 감사 가능하며, 머신러닝 인프라가 필요 없다. LLM 기반 계획기에 투자하기 전에 규칙 기반 적응형 라우팅을 1차 계획 메커니즘으로 구현할 것을 권고한다.

\textbf{시사점 3: LLM 기반 계획이 상당한 비용으로 순위 품질 이점을 제공한다.} GPT-5 계획기의 MRR(+2.0\%)과 nDCG(+1.2\%) 이점은 처리량이나 비용 제약보다 답변 정밀도가 더 중요한 고가치, 저용량 검색 응용(임원 의사결정 지원, 컴플라이언스 감사, 법적 연구)에서 \$0.016/질의 비용을 정당화할 수 있다. 고용량 기업 검색 워크로드에서는 비용-편익 비율이 규칙 계획기에 유리하다.

\textbf{시사점 4: 예외 처리 및 시간적 질의에는 전문화된 해결책이 필요하다.} 어떤 평가 시스템도 예외 질의(Recall@10 $\leq{0.412}$)나 시간적 질의(Dense: 0.000, 최선: 0.200)에서 허용 가능한 성능(Recall@10 $>{0.5}$)을 달성하지 못한다. 이러한 실패 유형은 매개변수적이 아닌 구조적이다 \citep{manning2008ir}. 예를 들어 날짜 필드 부스팅이 있는 시간적 인식 희소 검색과 정책 조항 쌍에 훈련된 예외 인식 밀집 모델과 같은 전용 처리가 필요하다.

\section{타당성 위협}
\label{sec:disc_threats_ko}

\subsection{내적 타당성}

\textbf{Oracle 레이블 한계.} Oracle 레이블의 대부분(약 87\%)이 규칙 계획기에 인코딩된 것과 동일한 결정론적 매핑에 의해 할당되었다. 이 구조적 정렬은 규칙 계획기의 SSA를 더 적대적인 평가와 비교하여 과장시킨다. GPT-5 오류의 44.7\%가 \texttt{oracle\_disagreement}로 분류된 것은 Oracle 레이블이 상당한 비율의 질의에서 유일하게 올바르지 않음을 나타낸다.

\textbf{청크 수준 정답.} 정답은 원본 문서의 축어적 증거 범위에서 파생된 필수 청크 ID 집합으로 구성된다. 섹션 인식 청크 분할은 경계에 걸쳐 증거를 분할할 수 있어, Recall@1 어려움을 인위적으로 증가시킬 수 있다.

\textbf{검색 지표 상한.} API 문서와 출장 정책은 복수 시스템에서 Recall@10$={1.000}$을 달성하여 $k{=}10$에서 이러한 범주에 대한 방법 간 판별이 불가능하다.

\subsection{외적 타당성}

\textbf{합성 데이터셋.} 합성 문서는 도메인 전문 용어, 상호 참조 패턴, 버전 특화 용어, 기업 코퍼스에서 흔한 편집 산물 등 실제 기업 문서의 전체 언어적 다양성을 포착하지 못할 수 있다.

\textbf{코퍼스 규모.} 120개 문서와 1,206개 청크로 SEKD는 소규모 코퍼스이다. BM25의 어휘 불일치와 밀집 검색의 임베딩 공간 포화가 더 두드러지는 대규모 환경(수만 개의 문서)에서 검색 방법이 다르게 동작할 수 있다.

\textbf{단일 임베딩 모델.} 밀집 검색에는 BGE-small-en-v1.5만 평가되었다. 더 크고 최신 모델(text-embedding-3-large, E5-large)은 특히 예외 및 시간적 질의에서 실질적으로 다른 순위를 생성할 수 있다.

\subsection{구성 타당성}

\textbf{계획 프록시로서의 SSA.} SSA는 oracle 동의를 측정하며, 사용자 대면 검색 품질을 측정하지 않는다. SSA 차이(GPT-5: 규칙 대비 -0.79 pp)와 검색 지표 차이(MRR: GPT-5에 유리 +1.05 pp) 사이의 약한 상관관계는 이러한 단절을 보여준다.
